\par\Needspace{18\baselineskip}\bigskip
\pdfbookmark[2]{Erster Abschnitt. Elliptische Bewegung.}{pallas-II-section1}
\begin{center}
{\large\bfseries Erster Abschnitt.\par}
\medskip\textit{Elliptische Bewegung.}
\end{center}
\begin{flushright}
\begin{minipage}{.94\textwidth}
\small\raggedright
. . . . . Ye other wandring fires that move\\
In mystic dance, not without song, resound\\
His praise, who out of darkness call’d up light.\\
\hspace*{2em}\textsc{Milton}, Paradise lost.
\end{minipage}
\end{flushright}
\pdfbookmark[3]{Darstellung, 1}{pallas-II-article1}
\begin{center}1.\end{center}
\noindent\hspace*{\parindent}Es seien \(x,y,z\) die Koordinaten eines Planeten in Bezug auf drei zueinander senkrechte Ebenen durch den Mittelpunkt der Sonne. Ferner sei \(t\) die Zeit, \(m\) die Summe der Massen der Sonne und des Planeten und \(r\) der Abstand des \SourceLeaf{447}Planeten von der Sonne. Sieht man von den Störungen ab, so erhält man
\[
\begin{array}{rl}
1)\quad&0=\dfrac{\mathrm{dd}x}{\mathrm{d}t^2}+\dfrac{mx}{r^3}\\[5pt]
2)\quad&0=\dfrac{\mathrm{dd}y}{\mathrm{d}t^2}+\dfrac{my}{r^3}\\[5pt]
3)\quad&0=\dfrac{\mathrm{dd}z}{\mathrm{d}t^2}+\dfrac{mz}{r^3}\mathpunct{\cdot}
\end{array}
\]
\par Setzt man folglich
\[
\begin{array}{rl}
4)\quad&\dfrac{y\mathrm{d}z-z\mathrm{d}y}{\mathrm{d}t}=A\\[5pt]
5)\quad&\dfrac{z\mathrm{d}x-x\mathrm{d}z}{\mathrm{d}t}=B\\[5pt]
6)\quad&\dfrac{x\mathrm{d}y-y\mathrm{d}x}{\mathrm{d}t}=C,
\end{array}
\]
so sind \(A,B,C\) konstante Größen, die daher vom Anfangszustand abhängen. Aus diesen Gleichungen erhält man
\[
\begin{array}{rl}7)\quad&Ax+By+Cz=0.\end{array}
\]
\par Die Bahn liegt also in einer Ebene, die durch den Mittelpunkt der Sonne geht. Ferner gilt
\[
\begin{gathered}
rr=xx+yy+zz\\
r\mathrm{d}r=x\mathrm{d}x+y\mathrm{d}y+z\mathrm{d}z\\
\begin{aligned}
r\mathrm{dd}r+\mathrm{d}r^2&=x\mathrm{dd}x+y\mathrm{dd}y+z\mathrm{dd}z+\mathrm{d}x^2+\mathrm{d}y^2+\mathrm{d}z^2\\
&=\mathrm{d}x^2+\mathrm{d}y^2+\mathrm{d}z^2-\frac{m\mathrm{d}t^2}{r}\mathpunct{\cdot}
\end{aligned}
\end{gathered}
\]
\par Nun liefern die Gleichungen 4, 5, 6
\[
\begin{aligned}
(AA+BB+CC)\mathrm{d}t^2
&=(xx+yy+zz)(\mathrm{d}x^2+\mathrm{d}y^2+\mathrm{d}z^2)\\
&\qquad{}-(x\mathrm{d}y+y\mathrm{d}y+z\mathrm{d}z)^2\\
&=rr(\mathrm{d}x^2+\mathrm{d}y^2+\mathrm{d}z^2)-rr\mathrm{d}r^2\\
&=r^3\mathrm{dd}r+mr\mathrm{d}t^2.
\end{aligned}
\]
\par Setzt man also
\[
\frac{AA+BB+CC}{m}=p,\qquad r-p=u,
\]
so erhält man \(0=\frac{\mathrm{dd}u}{\mathrm{d}t^2}+\frac{u}{r^3}\mathpunct{\raisebox{.45ex}{,}}\) eine Gleichung, die den ersten drei entspricht. Man sieht also, dass, wenn man
\[
\begin{array}{rl}
8)\quad&\dfrac{x\mathrm{d}u-u\mathrm{d}x}{\mathrm{d}t}=F\\[5pt]
9)\quad&\dfrac{y\mathrm{d}u-u\mathrm{d}y}{\mathrm{d}t}=G\\[5pt]
10)\quad&\dfrac{z\mathrm{d}u-u\mathrm{d}z}{\mathrm{d}t}=H,
\end{array}
\]
\SourceLeaf{448}setzt, \(F,G,H\) ebenfalls konstante Größen sind, die vom Anfangszustand abhängen. Verbindet man diese Formeln mit 4, 5, 6, so erhält man
\[
\begin{array}{rl}
11)\quad&Gz-Hy=A(r-p)\\
12)\quad&Hx-Fz=B(r-p)\\
13)\quad&Fy-Gx=C(r-p).
\end{array}
\]
\par Aus diesen Gleichungen folgt leicht, dass die Projektionen der Bahn und die Bahn selbst Kurven zweiter Ordnung sind, dass die Bahn den Halbparameter \(p\) hat und dass der Mittelpunkt der Sonne einen ihrer Brennpunkte einnimmt. Ferner liefern sie unmittelbar eine Bedingungsgleichung zwischen den sechs Konstanten \(A,B,C,F,G,H\)
\[
\begin{array}{rl}14)\quad&AF+BG+CH=0.\end{array}
\]
\par Ebenso erhält man aus den Gleichungen 8, 9, 10
\[
\begin{aligned}
&(FF+GG+HH)\mathrm{d}t^2\\
&=(xx+yy+zz)\mathrm{d}u^2-2(x\mathrm{d}x+y\mathrm{d}y+z\mathrm{d}z)u\mathrm{d}u\\
&\qquad{}+(\mathrm{d}x^2+\mathrm{d}y^2+\mathrm{d}z^2)uu\\
&=rr\mathrm{d}r^2-2r\mathrm{d}r\mathpunct{.}(r-p)\mathrm{d}r+(\mathrm{d}x^2+\mathrm{d}y^2+\mathrm{d}z^2)(r-p)^2\\
&=(AA+BB+CC)\left(1-\frac{2p}{r}\right)\mathrm{d}t^2+(\mathrm{d}x^2+\mathrm{d}y^2+\mathrm{d}z^2)pp.
\end{aligned}
\]
\par Setzt man also
\[
\begin{array}{rl}15)\quad&\dfrac{FF+GG+HH}{AA+BB+CC}=1-\dfrac{p}{a}\mathpunct{\raisebox{.45ex}{,}}\end{array}
\]
so erhält man
\[
\begin{array}{rl}16)\quad&0=\dfrac{\mathrm{d}x^2+\mathrm{d}y^2+\mathrm{d}z^2}{\mathrm{d}t^2}-\dfrac{2m}{r}+\dfrac{m}{a}\mathpunct{\cdot}\end{array}
\]
\par Man hätte dieselbe Gleichung erhalten können, indem man die Gleichungen 1, 2, 3 mit \(2\mathrm{d}x,\ 2\mathrm{d}y,\ 2\mathrm{d}z\) multipliziert und die Summe integriert. Die vorstehende Herleitung hat den Vorteil, den Zusammenhang zwischen der Konstanten, die diese Integration eingeführt hätte, und den schon gefundenen Konstanten zu zeigen.
\par
